\SourceLeaf{231}%
\Article{6}
\indent When the direction of the normal to the curved surface is transferred to the sphere, a definite point on the sphere corresponds to each definite point on the curved surface. In the same way, any line or figure on the curved surface will be represented by a corresponding line or figure on the sphere. In comparing two figures that correspond to one another in this way, one of which will be, as it were, an image of the other, we must consider two aspects. In one we consider quantity alone. In the other we set aside quantitative relations and consider position alone.

The first aspect will form the basis of certain notions that it seems useful to introduce into the theory of curved surfaces. To every part of a curved surface enclosed within definite boundaries we shall assign a \textit{total} or \textit{integral curvature}, expressed by the area of the corresponding figure on the sphere. We must carefully distinguish this integral curvature from a kind of specific curvature, which we shall call the \textit{measure of curvature}. The latter refers to a \textit{point} of the surface and denotes the quotient obtained by dividing the integral curvature of the surface element adjoining the point by the area of that element itself. It therefore gives the ratio of the corresponding infinitely small areas on the curved surface and on the sphere. We hope that what we explain later will amply confirm the usefulness of these innovations. As for the terminology, we thought it especially important to prevent any ambiguity. We therefore thought it unsuitable to follow strictly the analogy of the terminology commonly accepted, though not approved by everyone, in the theory of plane curves. By that analogy we would have had to call the measure of curvature simply curvature, and the integral curvature amplitude. But why should we not be free in our choice of words, provided that the subject has substance and the wording leaves no room for misinterpretation?

The position of a figure on the sphere can be similar to that of the corresponding figure on the curved surface or opposite (reversed). The former is the case when two lines that leave the same point on the curved surface in different directions that are not opposite are represented on the sphere by lines similarly placed, for example when the image of the line on the right is itself on the right. The latter is the case when the reverse holds. We shall distinguish these two cases by the positive or negative \textit{sign} of the measure of curvature. Evidently this distinction applies only insofar as, on each of the two surfaces, %
